\section{Introduction}
\label{sec:intro}

Threshold cryptography solves a fundamental problem in digital security: how
to distribute trust so that no single party---no single device, no single
employee, no single server---can unilaterally access, authorize, or destroy
a protected resource. The theoretical foundations are decades old. Shamir's
secret sharing \cite{shamir1979share} dates to 1979. Threshold signature
schemes have been refined through Gennaro and Goldfeder \cite{gennaro2018fast},
FROST \cite{komlo2020frost}, and CGGMP21 \cite{canetti2021cggmp}. Fully
homomorphic encryption enables threshold policy evaluation without
decryption \cite{gentry2009fhe}.

Despite this maturity, threshold cryptography has near-zero penetration in
consumer software. The handful of deployments---multisig wallets in
cryptocurrency, HSM quorum policies in enterprise key management---are
used exclusively by technical specialists who understand the underlying
cryptographic model. No mainstream application presents threshold
operations as a natural product experience.

The gap is not technical. It is a design failure. Threshold cryptography
is presented to users in its own terms: ``distribute $n$ shares, assemble
$t$ shares, reconstruct the secret.'' These terms mean nothing to a user
who wants to share a document with their team, recover their account after
losing their phone, or require two signatures on a large payment. The
cryptographic ceremony is visible, confusing, and intimidating.

This paper closes the gap by defining Threshold UX: a mapping from
threshold cryptographic primitives to familiar product interaction patterns
on Lux Network. The key insight is that threshold operations already have
natural product analogues:

\begin{itemize}
    \item A $t$-of-$n$ threshold signature is an \emph{approval flow}: ``2 of 3 team members must approve this action.''
    \item Shamir secret sharing is \emph{social recovery}: ``Any 3 of your 5 trusted contacts can help you regain access.''
    \item Threshold access control is a \emph{shared workspace}: ``Everyone on the team can read; edits require approval.''
    \item Threshold spend authorization is \emph{treasury management}: ``Large payments need two signers.''
    \item Threshold reveal is \emph{regulated access}: ``A compliance officer must co-approve before viewing PII.''
\end{itemize}

The user never sees shares, quorums, or reconstruction. They see approvals,
guardians, team permissions, and spending limits. The cryptography is real.
The ceremony is gone.

Section~\ref{sec:primitives} defines the threshold primitives available on
Lux. Section~\ref{sec:approval} maps these to approval flows.
Section~\ref{sec:recovery} addresses social recovery.
Section~\ref{sec:workspaces} covers shared workspaces.
Section~\ref{sec:treasury} formalizes treasury management.
Section~\ref{sec:regulated} describes regulated access.
Section~\ref{sec:progressive} introduces progressive security.
Section~\ref{sec:ux} defines UX interaction patterns.
Section~\ref{sec:security} proves security properties.
Section~\ref{sec:conclusion} concludes.

